Modern multistage decomposition is the closest computational comparison. \citet{Girardeau2015} prove almost-sure convergence for a class of sampling-based nested decomposition algorithms with general convex stage costs. \citet{Dowson2020} formulate policy graphs and extend stochastic dual dynamic programming to a class including convex infinite-horizon problems with continuous states and controls. \citet{FullnerRebennack2025} survey the assumptions, variants, and enhancements of stochastic dual dynamic programming. \citet{Lan2023} analyze its iteration complexity. These results are substantially broader than our tractable contract subclass. Merely choosing a shared design variable, adding a promise state, or normalizing a continuation multiplier is not a distinction from that literature.

The residual conclusion here is different: scalar continuation caps act by truncating a monotone payment response, each cap contributes at most one new breakpoint, and a single finite union suffices across all vertices of a recombining public input graph. The construction generates those thresholds before the root optimum is known. It gives an arithmetic storage and construction bound, and identifies the changing information record needed for execution. It does not assert that the number of generic Benders cuts or the bit cost of exact arithmetic is uniformly small.

Separable resource allocation is also a close comparison, not an omitted alternative. \citet{Vidal2016} develop decomposition for nested resource allocation; \citet{SchootUiterkamp2021} give a fast algorithm for quadratic allocation with nested constraints; \citet{Wu2021} develop an exact combinatorial method for separable convex allocation with nested bound constraints. Their resource-allocation algorithms make monotone prices and threshold operations established tools. We do not claim a new water-filling principle. The result proved here measures complexity in the compact stochastic public graph, retains different inherited prices when histories recombine, and derives the exact memory-value frontier for the stated renewal institution. Applying an algorithm to an already unfolded resource-allocation instance does not by itself give that compact-input bound. The full-tree and graph-sized descriptions are explicitly distinguished in the proof.

Nonanticipativity and policy aggregation \citep{RockafellarWets1991}, shared-design Benders masters \citep{Benders1962}, and stochastic dual value approximations \citep{PereiraPinto1991} remain the appropriate descriptions of the general representation in the print appendices. The finite-price theorem supplies a constructive stronger conclusion for a specified subclass, not a renamed version of those general objects. The accompanying result-by-result comparison records both the inherited principles and the precise additional claim being submitted for scrutiny.
